\section{Definición y función general}

El citoesqueleto es una red dinámica de filamentos proteicos que proporciona soporte estructural, organización espacial y movimiento a la célula. Está compuesto por tres tipos principales de filamentos.

\section{Tipos de filamentos}

\subsection{Filamentos intermedios}

\begin{itemize}
    \item Diámetro: 8-12 nm
    \item Composición: Proteínas fibrosas (queratina, vimentina, láminas)
    \item Función principal: Resistencia mecánica y soporte estructural
    \item Características: Mayor resistencia a la tensión
    \item Localización: Núcleo y citoplasma
\end{itemize}

\subsection{Microtúbulos}

\begin{itemize}
    \item Diámetro: 25 nm
    \item Composición: Tubulina (α y β)
    \item Función principal: Transporte intracelular y división celular
    \item Características: Estructura polar, dinámicos
    \item Localización: Huso mitótico, cilios, flagelos
    \item Asociados con: Quinesinas y dineínas (proteínas motoras)
\end{itemize}

\subsection{Filamentos de actina (microfilamentos)}

\begin{itemize}
    \item Diámetro: 5-7 nm
    \item Composición: Actina G (monómeros) y Actina F (polímeros)
    \item Función principal: Movimiento celular y contracción
    \item Características: Más delgados y flexibles
    \item Localización: Cortex celular, microvellosidades
    \item Asociados con: Miosina (contracción muscular)
\end{itemize}

\section{Funciones del citoesqueleto}

\begin{enumerate}
    \item \textbf{Soporte estructural:} Mantiene la forma de la célula
    \item \textbf{Organización intracelular:} Posiciona organelos
    \item \textbf{Transporte:} Movimiento de vesículas y organelos
    \item \textbf{Movimiento celular:} Migración y cambios de forma
    \item \textbf{División celular:} Formación del huso mitótico
    \item \textbf{Contracción muscular:} Interacción actina-miosina
\end{enumerate}

\section{Proteínas asociadas}

\begin{itemize}
    \item Miosina: Proteína motora asociada a actina
    \item Quinesina: Transporta cargas hacia el extremo (+) de microtúbulos
    \item Dineína: Transporta cargas hacia el extremo (-) de microtúbulos
\end{itemize}